\section{Phase 2 Implementation Roadmap}
\label{sec:future:roadmap}

Phase~2 of this thesis is the implementation of the dashboard
specified in Chapter~\ref{chapter:proposed}. The work is
organised into nine phases, P1--P9, each producing a working
demoable artefact. The phases are summarised in
Table~\ref{tab:future:roadmap}; their critical path is
P1~$\rightarrow$~P2 and the remaining phases are largely
parallelisable. Per Câu~5 of the grading
rubric~\cite{rubric}, the plan covers content, time, personnel,
and expected outcomes; full per-phase deliverables and
acceptance criteria are tabulated in
Appendix~\ref{appendix:plan}.

\begin{table}[H]
\centering
\caption{Phase 2 implementation phases, deliverables, and
estimated effort.}
\label{tab:future:roadmap}
\footnotesize
\renewcommand{\arraystretch}{1.20}
\begin{tabularx}{\linewidth}{p{0.9cm}p{4.2cm}Xp{1.4cm}}
\hline
\textbf{Ph.} & \textbf{Deliverable} & \textbf{Outcome} & \textbf{Effort} \\
\hline
P1 & Warehouse Layers 2--4 build
   & \texttt{dim\_*}, \texttt{fact\_*}, three \texttt{agg\_*}
     tables; build CLI; validation notebook 05
   & 2~days \\
P2 & End-to-end Card~2 (Demand)
   & Five charts render; linked brushing across all five; A3
     entity panel for routes; cold-load $<$~3~s; warm render
     $<$~50~ms
   & 3~days \\
P3 & Spatial card + pydeck
   & Stop Demand Explorer, hex density grid, network overview;
     \texttt{agg\_stop\_hour} populated
   & 2~days \\
P4 & School flagship card
   & Schools $\rightarrow$ Routes treemap, school-routeness
     classifier, equity map (district stubs)
   & 2~days \\
P5 & Financial + Network Health cards
   & P\&L quadrant, operator scorecard, payment adoption,
     QoS-1/2/3 scorecard, anomaly threshold-$T$ slider
   & 3~days \\
P6 & Routes card with Marey
   & Discovery (UMAP, archetype) and Detail modes; embedded
     Marey and GroupSet HTML components
   & 3~days \\
P7 & Hub + entity-panel polish
   & Editorial hub HTML embedded; A3 panel for all six entity
     types; URL state; export
   & 2~days \\
P8 & Stories scrollytelling
   & Four scenes-driven stories with Scrollama; scene JSON
     authoring
   & 2~days \\
P9 & Final polish + evaluation
   & Verified school geocoding; full 50-day GPS loaded if
     available; perf tuning; screen recording; initial user
     evaluation
   & 2~days \\
\hline
\multicolumn{3}{r}{\textbf{Total estimated:}} & 21~days \\
\hline
\end{tabularx}
\end{table}

\begin{figure}[H]
\centering
\imageplaceholder{Diagram (Gantt chart)}{A horizontal Gantt
chart with nine bars labelled P1--P9, the x-axis covering days
0--21. P1 (days 0--2, blue), P2 (days 2--5, blue), P3 (days
5--7, green; can start after P1+P2), P4 (days 5--7, gold; can
parallel with P3), P5 (days 7--10, red), P6 (days 10--13,
purple), P7 (days 13--15, cyan), P8 (days 15--17, orange), P9
(days 17--21, grey). Dependency arrows: P1 $\rightarrow$ P2;
P2 $\rightarrow$ P3, P4; P3, P4 $\rightarrow$ P5;
P5 $\rightarrow$ P6; P6 $\rightarrow$ P7; P7 $\rightarrow$ P8;
P8 $\rightarrow$ P9.}{Draw in a tool such as
\texttt{draw.io}, \texttt{Lucidchart}, or \texttt{Excalidraw};
export as PNG at $\sim$300~DPI. Save as
\texttt{images/diagrams/phase2\_gantt.png}.}
\caption{Phase 2 Gantt: nine phases over approximately 21
days. P3 and P4 can run in parallel after P1+P2 complete.}
\label{fig:future:gantt}
\end{figure}

\paragraph{Content.}
The plan content is the nine phases above. Each phase produces
one or two artefacts that the next phase depends on; the
Appendix~\ref{appendix:plan} table records the exact acceptance
criterion per phase (e.g., for P1: the build CLI runs in
$<$~30~s on the local sample and emits eight aggregate tables).

\paragraph{Time.}
Total estimated effort is twenty-one working days at one
developer. The critical path is P1~$\rightarrow$~P2 (five
working days); the remaining sixteen days are parallelisable in
P3+P4 (the spatial and the school cards run independently
because their aggregate tables \texttt{agg\_stop\_hour} and
\texttt{agg\_school\_route} are independent) and otherwise
sequential.

\paragraph{Personnel.}
Phase~2 is single-author. Time tracking is per-phase rather
than per-task because every phase is small enough for one
developer to hold the full context. The thesis advisor reviews
at the end of P2 (first end-to-end card) and at the end of P6
(all cards rendering) before sign-off.

\paragraph{Expected outcomes.}
At the end of Phase~2, the dashboard renders thirty out of
thirty-one charts at sub-fifty-millisecond warm latency;
supports linked brushing within every card and the universal
entity-detail panel (A3 of
Section~\ref{sec:proposed:invariants}) on every clickable
element; persists state to a shareable URL; and exports every
chart as PNG/PDF/CSV with a single click. The final chart, the
Marey time--distance diagram, depends on the full GPS pull
(Section~\ref{sec:future:gps}) and is delivered at the start
of P9 if the data is available, or deferred to a follow-on
phase if not.

\paragraph{Risk and mitigation.}
Four risks dominate the schedule. (1)~Streamlit performance
falling below the budget
(Section~\ref{sec:proposed:problem}); the mitigation is to fall
back to embedded HTML components for hot paths, which is
already the case for three charts. (2)~The full GPS pull from
Kaggle failing or taking longer than expected; the mitigation
is to demo with the local one-day sample and document the
seven charts that remain at SHOULD/WON'T tier until the pull
completes. (3)~School geocoding under target before P9
deadline; the mitigation is to ship the dashboard with
district-level approximation and the
``verified~\checkmark'' badge described in
Section~\ref{sec:proposed:warehouse_local}.
(4)~Linked-brushing latency feeling sluggish on Streamlit's
server-side rerun model; the mitigation is to pre-compute the
highlight sets server-side so only the chart's colour column
is recomputed on rerun, not the data itself.
